\chapter{Comparative Analysis}

This chapter provides a cross-sample comparison of the three Malbus APK variants analysed in the previous chapters: \textbf{Malbus v1.0.3}, \textbf{Malbus v3.3.7}, and \textbf{Malbus v3.6.5}.
The goal is to identify common patterns, progressive evolution of malicious capabilities, and differences in evasion sophistication across the family.

\section{VirusTotal Detection}

All three samples were submitted to VirusTotal and flagged as malicious by a comparable number of security engines, consistently classified as Android trojans of the downloader family.
The dominant threat label across all samples is \texttt{trojan.andr/apkdownloader}, and the Tencent engine correctly identified each sample by referencing its specific APK package name.

\begin{table}[H]
\centering
\begin{tabular}{|l|c|c|c|}
\hline
\textbf{Metric} & \textbf{v1.0.3} & \textbf{v3.3.7} & \textbf{v3.6.5} \\
\hline
Malicious detections & 40/75 & 40/75 & 38/75 \\
Reputation score & $-1$ & $-2$ & $0$ \\
Dominant threat label & trojan & trojan/downloader & trojan/downloader \\
Total submissions & 377 & 354 & -- \\
Alternative filenames & Yes & Yes & Yes \\
\hline
\end{tabular}
\caption{VirusTotal detection summary across the three APK samples}
\end{table}

Notably, well-known vendors such as Microsoft Defender, McAfee, and Malwarebytes failed to detect any of the three samples as malicious, highlighting a significant gap in coverage for this malware family.

\section{Permissions and Attack Surface}

All three APKs share a largely identical permission set, which is consistent with them belonging to the same malware family.
The most security-relevant permissions present in all samples are:

\begin{itemize}
    \item \texttt{ACCESS\_FINE\_LOCATION} and \texttt{ACCESS\_COARSE\_LOCATION} --- device tracking via GPS and cell towers.
    \item \texttt{READ\_EXTERNAL\_STORAGE} / \texttt{WRITE\_EXTERNAL\_STORAGE} --- access to user files on external memory.
    \item \texttt{REQUEST\_INSTALL\_PACKAGES} --- silent installation of additional APKs, a key dropper indicator.
    \item \texttt{INTERNET} --- unrestricted network access, used for remote configuration and payload delivery.
    \item GCM/FCM push messaging --- allows remote triggering of behaviour.
\end{itemize}

The main difference is that v1.0.3 additionally declares \texttt{WRITE\_HISTORY\_BOOKMARKS} and \texttt{READ\_GSERVICES}, which are absent in the later versions, suggesting a minor refinement of the permission scope in subsequent releases.

\section{Behavioural Tags and Evasion Capabilities}

A clear evolution in anti-analysis sophistication is observable across the three versions.
While v1.0.3 already employs obfuscation and Java Reflection, versions v3.3.7 and v3.6.5 introduce explicit VM and debugger evasion mechanisms, as confirmed by the VirusTotal behavioural tags.

\begin{table}[H]
\centering
\begin{tabular}{|l|c|c|c|}
\hline
\textbf{Behavioural Tag} & \textbf{v1.0.3} & \textbf{v3.3.7} & \textbf{v3.6.5} \\
\hline
obfuscated & \checkmark & \checkmark & \checkmark \\
reflection & \checkmark & \checkmark & \checkmark \\
checks-gps & \checkmark & \checkmark & \checkmark \\
telephony & \checkmark & \checkmark & \checkmark \\
detect-debug-environment & $\times$ & \checkmark & \checkmark \\
checks-cpu-name & $\times$ & \checkmark & \checkmark \\
sets-process-name & $\times$ & \checkmark & \checkmark \\
\hline
\end{tabular}
\caption{VirusTotal behavioural tags across the three APK samples}
\end{table}

This progression indicates that the malware author actively improved the evasion layer between v1.0.3 and the later variants, adding sandbox detection and process name manipulation to harden the payload against dynamic analysis environments.

\section{Static Analysis: Common Architecture}

The static analysis of all three samples reveals an identical internal architecture, confirming they are regional variants of the same malware family.
The core classes (\texttt{MainActivity}, \texttt{BaseMainActivity}, \texttt{EnvConstant}, \texttt{BusApiWeb}, \texttt{AppInfoActivity}) are present and functionally equivalent across all three APKs, differing only in the regional bus endpoint URLs and the specific package names (\texttt{cwbus}, \texttt{gjbus}, \texttt{jjbus}).

The following structural elements are common to all three samples:

\begin{itemize}
    \item Loading of a native library via \texttt{System.loadLibrary("Audio3.0")} in \texttt{BaseMainActivity}.
    \item Native methods \texttt{startUpdate()} and \texttt{updateApplication()} consistent with dropper-like update logic.
    \item Remote environment loading mechanism that allows runtime behaviour modification without a new APK release.
    \item Hardcoded advertising and analytics keys in \texttt{EnvConstant} (AdMob, AdLib, Facebook, Google Analytics).
    \item Unencrypted HTTP communication in \texttt{BusApiWeb}, exposing data to man-in-the-middle attacks.
    \item A cosmetic \texttt{AppInfoActivity} displaying a legitimate bus service URL to disguise the application's true nature.
\end{itemize}

\section{MobSF Static Analysis Comparison}

MobSF confirms consistent and severe security issues across all three samples.
All APKs share the same certificate signing profile: SHA-1 with RSA, v1 scheme only, issued by CN=Myoung-Hun Lee --- making all three vulnerable to the Janus attack on Android 5.0--8.0.

\begin{table}[H]
\centering
\begin{tabular}{|l|c|c|c|}
\hline
\textbf{Finding} & \textbf{v1.0.3} & \textbf{v3.3.7} & \textbf{v3.6.5} \\
\hline
Janus vulnerability (v1 signing only) & \checkmark & \checkmark & \checkmark \\
SHA-1 weak hashing & \checkmark & \checkmark & \checkmark \\
Low minSdk (14) & \checkmark & \checkmark & \checkmark \\
Task Hijacking / StrandHogg & \checkmark & \checkmark & \checkmark \\
World-readable file / shared preference & \checkmark & \checkmark & \checkmark \\
Hardcoded secrets & \checkmark & \checkmark & \checkmark \\
Privacy trackers ($>$8) & \checkmark & \checkmark & \checkmark \\
INSTALL\_PACKAGES external permission & \checkmark & \checkmark & \checkmark \\
High severity findings (count) & 9 & 9 & 9 \\
\hline
\end{tabular}
\caption{MobSF high-severity findings across the three APK samples}
\end{table}

The identical number of high-severity findings (9) across all three samples further confirms the shared codebase and the lack of any remediation effort between releases.

\section{YARA Rule Matches}

All three samples matched the \texttt{IcedID} structural YARA rule from the \texttt{CAPEv2} repository (author: \texttt{kevoreilly/threathive}).
Although IcedID is historically associated with a Windows banking trojan, the matched rule targets dropper and loader architectures, not the specific malware family.
This consistent structural match across all three variants reinforces the hypothesis that the Malbus family was deliberately engineered with an advanced payload-delivery architecture, likely to withstand analysis and deliver second-stage payloads on demand.

\section{Overall Family Assessment}

The comparative analysis confirms that Malbus v1.0.3, v3.3.7, and v3.6.5 are three regional variants of the same trojanised application family, sharing an identical malicious core while adapting the legitimate-facing functionality to different South Korean cities (Changwon, Gwangju, Jeonju).

The key evolution across versions is the progressive addition of anti-analysis and VM evasion capabilities, which are absent in v1.0.3 but fully present in v3.3.7 and v3.6.5.
In all other dimensions --- permissions, architecture, signing, MobSF findings, advertising SDKs, and dropper logic --- the three samples are functionally identical, indicating a mature and deliberately maintained malware family rather than opportunistic or independent variants.
